%% ============================================================
%% 4. EXPERIMENTS
%% ============================================================
\section{Experiments}
\label{sec:experiments}

\subsection{Datasets}
\label{sec:datasets}

\paragraph{OakInk2~\cite{zhan2024oakink2} (Primary Benchmark).}
A bimanual hand-object manipulation dataset captured with SMPL-X~\cite{pavlakos2019expressive} body tracking and multi-object 6DoF annotations.
We extract 963 super-segments from 627 complex sequences by merging adjacent primitives with temporal gaps ${\leq}\,300$ frames.
Each segment undergoes boundary expansion (${\pm}60$ frames), downsampling from 30 to 10~FPS, and uniform subsampling to $T{=}200$ frames, yielding 534D dynamic features (267D positions $+$ 267D velocities) with up to 4 object slots.
Per-entity state labels $[T, 7]$ are computed via Schmitt-trigger contact detection ($\delta_\text{on}{=}1.5$~cm, $\delta_\text{off}{=}3.0$~cm) with majority-vote smoothing.
Train/test split: 770/193 (80/20).

\paragraph{GRAB~\cite{taheri2020grab}.}
A single-hand grasping dataset with 1{,}335 sequences across 10 subjects and 51 objects.
Raw SMPL-X parameters at 120~FPS are downsampled to 30~FPS.
State labels are derived from ground-truth per-vertex contact annotations on all 10{,}475 SMPL-X vertices.
Target length: 300 frames.
Subject-unseen split: train on s1--s8 (1{,}065), test on s9--s10 (270).

\paragraph{OMOMO~\cite{li2023omomo}.}
A full-body manipulation dataset focusing on large objects (furniture, appliances), with 4{,}294 sequences at 30~FPS from 17 subjects.
Features are extracted from the InterAct~\cite{xu2025interact} unified 962D representation by selecting the compatible 534D subset.
Train/test: 3{,}864/430.


\subsection{Baselines}
\label{sec:baselines}

We compare against two representative methods that address different aspects of the problem:

\paragraph{InterAct HOIDiff~\cite{xu2025interact}.}
A text-conditioned diffusion model for single-segment human-object interaction generation.
HOIDiff employs a dual transformer backbone with PointNet2/BPS object encoding, trained on the InterAct consolidated dataset (11K$+$ sequences).
For fair comparison on OakInk2, we train HOIDiff from scratch on our OakInk2 primitive data for 20K steps following the same preprocessing.
HOIDiff generates 962D motions (markers $+$ velocities $+$ contact features) but lacks per-entity state conditioning and long-horizon autoregressive capability.

\paragraph{PriorMDM DoubleTake~\cite{shafir2024priormdm}.}
A long-horizon body motion composition method that chains multiple motion segments via a two-pass ``handshake + blend'' refinement strategy.
DoubleTake uses a pre-trained MDM~\cite{tevet2023human} backbone on HumanML3D~\cite{guo2022generating} (14K sequences, 263D body-only representation).
It generates 22 body joints without hand articulation or object interaction---we include it as a long-horizon body motion baseline.


\subsection{Evaluation Metrics}
\label{sec:metrics}

Following the InterAct benchmark protocol~\cite{xu2025interact}, we evaluate on three categories:

\paragraph{Generation quality.}
\textbf{FID}~\cite{heusel2017gans}: Fr\'{e}chet Inception Distance computed using InterAct's pre-trained motion encoder (615D input: 231D markers $+$ 384D BPS-encoded object, 256D latent).
\textbf{Matching Score}: text-motion alignment via L2 distance between encoded text and motion embeddings; lower is better.
\textbf{Diversity}: pairwise L2 distance variance in motion embedding space.

\paragraph{Reconstruction quality.}
\textbf{Marker MSE}: mean squared error of 77 marker positions vs.\ ground truth.
\textbf{Object Pose Error}: position error (mm) and geodesic rotation error (degrees) for the primary object.
\textbf{Jitter}: third-order temporal derivative magnitude, measuring motion smoothness.

\paragraph{Interaction quality.}
\textbf{Contact F1}: precision and recall of hand-object contact at a 5~cm distance threshold, averaged over both hands.


\subsection{Main Results on OakInk2}
\label{sec:main_results}

\Cref{tab:main_oakink2} reports the primary comparison on the OakInk2 test set (193 sequences).
All metrics are computed with InterAct's pre-trained evaluator for generation quality, and directly from generated/ground-truth marker positions for reconstruction and interaction metrics.

\begin{table*}[t]
  \caption{\textbf{Main results on OakInk2.} Generation quality metrics (FID, Matching, Diversity) are computed using InterAct's pre-trained evaluator; reconstruction and interaction metrics are computed directly. Best results in \textbf{bold}. $\downarrow$: lower is better; $\uparrow$: higher is better. $^\dagger$: body-only (no hands/objects). $^\ddagger$: trained 20K steps on OakInk2.}
  \label{tab:main_oakink2}
  \centering
  \small
  \begin{tabular}{@{}l|ccc|cccc|c@{}}
    \toprule
    \multirow{2}{*}{Method} & \multicolumn{3}{c|}{Generation Quality} & \multicolumn{4}{c|}{Reconstruction} & Interaction \\
    & FID $\downarrow$ & Matching $\downarrow$ & Diversity $\uparrow$ & Marker MSE $\downarrow$ & Obj Pos (mm) $\downarrow$ & Jitter $\downarrow$ & Vel MSE $\downarrow$ & Contact F1 $\uparrow$ \\
    \midrule
    PriorMDM DoubleTake$^\dagger$ & 157.0 & 13.836 & 5.302 & --- & --- & 0.002$^*$ & --- & --- \\
    InterAct HOIDiff$^\ddagger$ & 2.821 & 14.229 & 1.514 & --- & --- & --- & --- & --- \\
    \midrule
    CRR-Flow (Ours) & 0.438 & 12.317 & 0.937 & \textbf{0.0140} & 223 & 0.0063 & \textbf{0.000056} & 0.036 \\
    \quad w/o AR (V4) & 0.075 & 12.312 & 1.407 & 0.0202 & \textbf{207} & \textbf{0.0051} & 0.000059 & 0.035 \\
    \quad w/o Vel+BPS (V3) & \textbf{0.058} & \textbf{12.279} & \textbf{1.401} & 0.0248 & 210 & \textbf{0.0051} & --- & \textbf{0.063} \\
    \bottomrule
  \end{tabular}
\end{table*}

\paragraph{Analysis.}
CRR-Flow achieves the lowest Marker MSE (0.014), reducing reconstruction error by 43\% compared to the position-only baseline (V3, 0.025).
In terms of FID, all our variants substantially outperform InterAct HOIDiff on OakInk2 (best: 0.058 vs.\ 2.821), indicating that our generated motions are distributionally closer to the ground truth.
The Matching Score similarly favors our approach (12.3 vs.\ 14.2), demonstrating better text-motion alignment.

We note that the Diversity values across all methods on OakInk2 (0.9--1.5) are significantly lower than those reported by InterAct on GRAB (7.7) or OMOMO (10.5).
This reflects the inherent characteristics of OakInk2: its laboratory manipulation tasks exhibit less kinematic variation than full-body activities such as furniture rearrangement.
The ground-truth Diversity on OakInk2 is 1.4--1.5, and our models closely match this distribution (V4: 1.41, V3: 1.40), confirming that the lower Diversity is not a model limitation.

PriorMDM DoubleTake generates body-only motions (22 joints, no hand articulation or objects) trained on a different domain (HumanML3D), precluding direct quantitative comparison on interaction metrics.
We include it as a qualitative long-horizon baseline in \cref{sec:long_horizon}.


\subsection{Ablation Study}
\label{sec:ablation}

We systematically ablate the three key components of CRR-Flow (\cref{tab:ablation}).
Starting from the full model (CRR-Flow = V4 AR+Release), we remove components incrementally.

\begin{table}[t]
  \caption{\textbf{Ablation study on OakInk2.} Each row removes one component from CRR-Flow. $\Delta$: relative change vs.\ full model in Marker MSE.}
  \label{tab:ablation}
  \centering
  \small
  \begin{tabular}{@{}lccc@{}}
    \toprule
    Configuration & Marker MSE $\downarrow$ & FID $\downarrow$ & $\Delta$ MSE \\
    \midrule
    \textbf{CRR-Flow} (full) & \textbf{0.0140} & 0.438 & --- \\
    \quad w/o Release state & 0.0141 & 0.438 & $+$0.7\% \\
    \quad w/o Autoregressive & 0.0202 & 0.075 & $+$44.3\% \\
    \quad w/o Vel $+$ BPS (V3) & 0.0248 & \textbf{0.058} & $+$77.1\% \\
    \bottomrule
  \end{tabular}
\end{table}

\paragraph{Velocity features and BPS geometry ($-$30.6\% MSE).}
Removing velocity features and BPS object geometry (V3 configuration, 267D input) increases Marker MSE from 0.020 to 0.025.
Velocity provides an explicit inductive bias for temporal dynamics: the model no longer needs to implicitly differentiate positions to predict frame-to-frame motion.
BPS geometry enables shape-aware manipulation---the object projection layer (1042D $\to$ 512D) compresses both dynamic pose (18D) and static shape (1024D) into a unified entity token.

\paragraph{Autoregressive relay ($-$30.7\% MSE).}
Removing the block-wise autoregressive strategy ($K{=}20$ history frames) increases Marker MSE from 0.014 to 0.020.
The history window provides a ``kinematic anchor'': rather than generating from pure noise, the model conditions on 2 seconds of established motion, grounding body position, velocity, and hand-object relationships.
This is our largest single-component improvement.

Interestingly, removing autoregressive conditioning \emph{improves} FID (0.075 vs.\ 0.438).
We attribute this to a distribution narrowing effect: the history window constrains generation toward the specific continuation implied by the conditioning, reducing the diversity of outputs.
Since FID measures distributional distance, the less constrained non-AR model produces a broader distribution that better matches the aggregate test set statistics, despite higher per-sample reconstruction error.

\paragraph{Release state ($+$0.7\% MSE).}
The Release state provides marginal quantitative improvement but is architecturally important for the autoregressive relay: it signals the model to decelerate and withdraw the hand before transitioning to the next primitive, enabling smoother inter-segment transitions in long-horizon generation.


\subsection{Long-Horizon Autoregressive Generation}
\label{sec:long_horizon}

A key capability of CRR-Flow is generating extended manipulation sequences through block-wise autoregressive relay.
We evaluate on 4 multi-segment chains from the OakInk2 test set, where consecutive super-segments from the same complex task are chained by relaying the last $K{=}20$ frames as history for the next block (\cref{tab:long_horizon}).

\begin{table}[t]
  \caption{\textbf{Long-horizon autoregressive generation.} CRR-Flow chains multiple segments via $K{=}20$ history relay. Single-segment Marker MSE is 0.014 for reference.}
  \label{tab:long_horizon}
  \centering
  \small
  \begin{tabular}{@{}lccc@{}}
    \toprule
    Task & Segments & Frames (Duration) & Marker MSE $\downarrow$ \\
    \midrule
    Rearrange plate & 4 & 740 (74s) & 0.0116 \\
    Scoop cheese  & 3 & 560 (56s) & 0.0185 \\
    Insert lightbulb & 3 & 560 (56s) & 0.0139 \\
    Take donut & 3 & 560 (56s) & 0.0167 \\
    \midrule
    \textbf{Average} & 3.25 & 605 (60.5s) & \textbf{0.0152} \\
    \bottomrule
  \end{tabular}
\end{table}

The average long-horizon Marker MSE (0.015) is comparable to single-segment quality (0.014), with only 8.6\% degradation despite generating up to 740 frames (74 seconds) of continuous manipulation.
This demonstrates that the history relay mechanism successfully preserves kinematic continuity across segment boundaries.
By comparison, PriorMDM DoubleTake can chain body motions of similar duration but cannot model hand-object contacts, object poses, or multi-entity state transitions---capabilities that are essential for manipulation tasks.
On a body-only jitter comparison (the only directly comparable metric), CRR-Flow achieves 2.2$\times$ lower body jitter (0.003 vs.\ 0.007) than PriorMDM, while additionally generating 27 hand markers and 4 object poses that PriorMDM cannot produce.
($^*$Note: PriorMDM jitter is computed on 22 SMPL joints; CRR-Flow jitter is on 50 SMPL-X body markers.)

The longest chain (\emph{Rearrange plate}, 4 segments) achieves the \emph{lowest} MSE (0.012), suggesting that simpler repeated actions benefit most from history conditioning.
The hardest chain (\emph{Scoop cheese}, MSE 0.019) involves complex bimanual coordination with tool use, which remains challenging for long-horizon generation.


\subsection{Cross-Dataset Evaluation}
\label{sec:cross_dataset}

To assess the generality of our entity-centric architecture, we train CRR-Flow on GRAB and OMOMO (\cref{tab:cross_dataset}).

\begin{table}[t]
  \caption{\textbf{Cross-dataset results.} CRR-Flow trained separately on each dataset. All use V4 architecture with autoregressive relay.}
  \label{tab:cross_dataset}
  \centering
  \small
  \begin{tabular}{@{}lccccc@{}}
    \toprule
    Dataset & Steps & M-MSE $\downarrow$ & Obj Pos $\downarrow$ & Jitter $\downarrow$ & F1 $\uparrow$ \\
    \midrule
    OakInk2 & 60K & \textbf{0.014} & \textbf{223} & \textbf{0.006} & 0.036 \\
    GRAB & 40K & 0.069 & 383 & 0.011 & \textbf{0.190} \\
    OMOMO & 40K & 0.386 & 1105 & 0.029 & 0.000 \\
    \bottomrule
  \end{tabular}
\end{table}

GRAB achieves the highest Contact F1 (0.19) because its ground-truth per-vertex contact labels provide precise state supervision, compared to OakInk2's distance-based Schmitt trigger approximation.
OMOMO presents the greatest challenge: its large furniture objects (suitcases, tables) require whole-body coordination patterns fundamentally different from the fine-grained hand manipulation in OakInk2 and GRAB.
The entity-centric tokenization and state conditioning framework generalizes across all three datasets without architectural changes, requiring only dataset-specific preprocessing.


\subsection{Implementation Details}
\label{sec:implementation}

CRR-Flow uses a DiT backbone with latent dimension $D{=}512$, 8 transformer layers, 8 attention heads, and FFN dimension 2048, totaling 41.2M trainable parameters (excluding frozen CLIP ViT-B/32).
We train with AdamW ($\text{lr}{=}10^{-4}$, weight decay $0.01$) for 200K steps with gradient clipping at 1.0 and batch size 32.
Flow matching uses independent coupling (ConditionalFlowMatcher, $\sigma{=}0$).
CLIP text embeddings are pre-computed and cached for memory efficiency.
Inference uses ODE integration with Euler's method (10 steps).
All experiments are conducted on NVIDIA RTX A6000 GPUs (48~GB).
